\documentclass[conference]{IEEEtran}
\IEEEoverridecommandlockouts

\usepackage{cite}
\usepackage{amsmath,amssymb,amsfonts}
\usepackage{graphicx}
\usepackage{textcomp}
\usepackage{xcolor}
\usepackage{booktabs}
\usepackage{multirow}
\usepackage{hyperref}
\usepackage{microtype}

\def\BibTeX{{\rm B\kern-.05em{\sc i\kern-.025em b}\kern-.08em
    T\kern-.1667em\lower.7ex\hbox{E}\kern-.125emX}}

\begin{document}

\title{Feature Sufficiency in Zero-Shot Antimicrobial Resistance Prediction:\\
When Molecular Fingerprints Outperform Graph Neural Networks}

\author{
  \IEEEauthorblockN{Anish Gupta}
  \IEEEauthorblockA{Department of Computer Science\\
    \textit{[Institution]}\\
    \texttt{[email]}}
}

\maketitle

\begin{abstract}
Zero-shot generalization---predicting which genes confer resistance to novel
antibiotic classes not seen during training---is a clinically critical task that
motivates the use of powerful inductive models such as heterogeneous graph
neural networks (GNNs). Prior work using CARD resistance data reported strong
zero-shot MRR scores ($0.775$), apparently validating complex architectures.
We show that these scores are inflated by \emph{data leakage}: the gene feature
vectors used contained explicit drug-class membership labels, directly encoding
the zero-shot answer. After replacing gene representations with ESM-2 protein
language model embeddings and drug representations with molecular fingerprints
(Morgan + MACCS + topological torsion), and adopting a rigorous evaluation
protocol (ranking each zero-shot drug class among all 46 drug classes rather
than a narrow 5-way subset), zero-shot performance drops dramatically.
Surprisingly, a two-layer feature-concatenation MLP with no graph structure
outperforms all tested GNNs---including a purpose-built heterogeneous GATConv
model---by a factor of \textbf{3.6$\times$} on the primary metric
($\text{ZS-all MRR} = 0.069$ vs.\ $0.019$; random = $0.022$). Ablation
studies reveal that \emph{both} gene (ESM-2) and drug (molecular fingerprint)
features are jointly necessary: removing either collapses ZS-all MRR to or
below random ($0.010$--$0.020$). The MLP learns pairwise gene-drug
compatibility from protein sequence and chemical structure jointly. Graph
topology provides no additional benefit and can actively harm
generalisation by over-smoothing training-class patterns onto held-out drugs. Our results serve as a methodological caution for the biomedical
GNN community: feature leakage is easy to overlook and drastically inflates
apparent zero-shot performance.
\end{abstract}

\begin{IEEEkeywords}
Antimicrobial resistance, zero-shot prediction, knowledge graph embedding, graph neural networks, data leakage, molecular fingerprints, protein language models
\end{IEEEkeywords}

% ─────────────────────────────────────────────────────────────────────────────
\section{Introduction}
\label{sec:intro}
% ─────────────────────────────────────────────────────────────────────────────

Antimicrobial resistance (AMR) is a growing global health crisis, responsible
for an estimated 1.27 million deaths annually and projected to surpass cancer
as a leading cause of mortality by 2050~\cite{murray2022global}. A key
challenge in computational AMR research is \emph{zero-shot} resistance
prediction: given a novel antibiotic class not encountered during model
training, predict which resistance genes confer resistance to it. This
capability would directly support surveillance of emerging resistance patterns
and the prioritisation of drug development targets.

Recent work has applied heterogeneous graph neural networks (GNNs) to the
CARD (Comprehensive Antibiotic Resistance Database)~\cite{alcock2023card},
building a knowledge graph with gene, drug-class, and mechanism nodes. These
methods have reported impressive zero-shot MRR scores (e.g., $0.775$), fuelling
interest in increasingly complex architectures~\cite{prior_work}. However, the
validity of these scores has not been critically examined.

In this paper we make the following contributions:
\begin{enumerate}
  \item We identify a \textbf{data leakage} flaw in prior gene feature vectors:
    the 46-dimensional drug-class membership component directly encodes which
    drug classes each gene resists, making zero-shot evaluation trivially easy.
    We quantify the inflation (reported $0.775$ vs.\ corrected $0.137$,
    a $5.7\times$ reduction) and provide an \emph{fmlp-leaky} ablation that
    reproduces the artefact.
  \item We introduce a \textbf{rigorous evaluation benchmark} using
    ESM-2~\cite{lin2023esm2} protein language model embeddings as gene
    features, RDKit molecular fingerprints as drug features, a 12-class
    zero-shot split, and an all-class ranking protocol (random = $0.022$
    vs.\ prior $0.477$ for 5-way ranking).
  \item Under this benchmark, a \textbf{simple Feature-MLP} achieves
    state-of-the-art zero-shot AMR performance
    ($\text{ZS-all MRR}=0.069 \pm 0.006$), outperforming all tested graph
    methods.
  \item \textbf{Ablations} reveal that \emph{both} ESM-2 gene embeddings and
    drug molecular fingerprints are jointly necessary for zero-shot transfer;
    removing either collapses ZS-all MRR to or below random. GNN processing
    of the resistance graph adds no benefit beyond flat feature concatenation.
  \item A \textbf{per-class analysis} reveals which antibiotic families are
    tractable (lincosamide, streptogramin A) and which remain fundamentally
    hard (fusidane, aminocoumarin), with performance correlating strongly with
    Tanimoto chemical similarity to training drug classes.
\end{enumerate}

\section{Background}
\label{sec:background}

\subsection{CARD and AMR Knowledge Graphs}
The Comprehensive Antibiotic Resistance Database (CARD 3.2.6)~\cite{alcock2023card}
catalogues 6,397 resistance gene models across 46 drug classes and
8 resistance mechanisms, with 12,593 gene-drug resistance associations. Its
structure naturally forms a heterogeneous knowledge graph with three node
types (gene, drug class, mechanism) and two edge types (gene-to-drug
resistance, gene-to-mechanism).

\subsection{Knowledge Graph Embedding for Link Prediction}
Classical embedding methods (TransE~\cite{bordes2013transe},
DistMult~\cite{yang2015distmult}) learn entity and relation embeddings from
observed triples. These methods are \emph{transductive}: they cannot generate
embeddings for unseen entities at test time. For zero-shot AMR prediction,
where the target drug class is unseen at training time, \emph{inductive}
methods that derive embeddings from node features are required.

R-GCN~\cite{schlichtkrull2018rgcn} extends GCN to heterogeneous graphs via
relation-type-specific convolutional filters, enabling inductive use when node
features are available. More recent heterogeneous GAT
models~\cite{yun2019heterogeneous} add attention mechanisms across edge types.

\subsection{Protein Language Models}
ESM-2~\cite{lin2023esm2} is a family of transformer-based protein language
models trained on 250 million protein sequences. The
\texttt{esm2\_t12\_35M\_UR50D} variant provides 480-dimensional embeddings via
mean-pooling over sequence token representations. Crucially, ESM-2 is trained
on sequences only, with no knowledge of resistance phenotypes---making its
embeddings leakage-free for AMR prediction tasks.

\subsection{Molecular Fingerprints}
Molecular fingerprints encode the structural and pharmacophoric properties of
small molecules as binary or count vectors. We concatenate three
complementary representations: Morgan circular fingerprints
(ECFP4, 2048 bits, radius 2)~\cite{rogers2010ecfp}, MACCS structural keys
(167 bits), and topological torsion fingerprints (1024 bits), yielding a
3,245-dimensional vector per drug class computed via
RDKit~\cite{landrum2016rdkit} from canonical SMILES strings. Tanimoto
similarity between these vectors preserves chemical neighbourhood structure.

\section{Methods}
\label{sec:methods}

\subsection{Graph Construction and Clean Features}
\label{sec:graph}

We build two versions of the CARD heterogeneous graph:

\textbf{Leaky graph} (reproducing prior work): gene features are 154-dimensional
vectors including a 46-bit drug-class membership one-hot, 8-bit mechanism
one-hot, and additional metadata. Drug features are 97-dimensional including a
46-bit one-hot class identity.

\textbf{Clean graph (BioMolAMR)}: gene features are 480-dimensional ESM-2
embeddings, L2-normalised. Drug features are 3,245-dimensional molecular
fingerprints. We additionally add drug-drug chemical similarity edges between
drug classes with Tanimoto similarity $>0.20$ (26 edges total), enabling
chemical neighbourhood information to propagate during drug encoding.

\subsubsection{Leakage Analysis}
The 46-bit drug-class membership component in the leaky gene features directly
encodes, for each gene, the set of drug classes to which it confers resistance.
For a held-out (zero-shot) drug class $d$, any gene $g$ resistant to $d$ has
the $d$-th component set to $1$, while non-resistant genes have it set to $0$.
A model receiving these features does not need to learn a meaningful
generalisation---it can simply read off the answer from the feature vector. We
demonstrate this quantitatively in \S\ref{sec:ablations}.

\subsection{Zero-Shot Evaluation Protocol}
\label{sec:splits}

We define a 12-class zero-shot split from the 46 CARD drug classes. The
12 held-out classes span diverse resistance mechanisms and were selected to
cover structural families not well-represented in the 34 training classes
(Table~\ref{tab:zs_classes}). This yields 9,053 training edges and 524
zero-shot test edges.

\textbf{Primary metric: ZS-all MRR.} For each zero-shot positive pair
$(g, d)$, we rank the true drug class $d$ against all 45 other drug classes
(negative sampling over all gene-drug pairs not in the positive set). This
gives a random baseline MRR of $1/46 \approx 0.022$---significantly harder
than the 5-way ZS-within MRR used in prior work (random = $0.477$).

\textbf{Secondary metric: ZS-within MRR.} Ranking among the 12 ZS classes
only (random = $1/12 \approx 0.083$). Reported for backward comparability.

\begin{table}[t]
\centering
\caption{12 Zero-Shot Drug Classes}
\label{tab:zs_classes}
\small
\begin{tabular}{lrr}
\toprule
Drug Class & \#Genes & Tan. to train \\
\midrule
Lincosamide antibiotic        & 108 & high \\
Streptogramin A antibiotic    &  53 & med  \\
Glycopeptide antibiotic       &  90 & med  \\
Rifamycin antibiotic          &  62 & low  \\
Pleuromutilin antibiotic      &  34 & low  \\
Phosphonic acid antibiotic    &  48 & low  \\
Isoniazid-like antibiotic     &  23 & low  \\
Oxazolidinone antibiotic      &  18 & low  \\
Glycylcycline                 &  35 & med  \\
Nitroimidazole antibiotic     &  18 & low  \\
Aminocoumarin antibiotic      &  24 & low  \\
Fusidane antibiotic           &  11 & low  \\
\bottomrule
\end{tabular}
\end{table}

\subsection{Models}
\label{sec:models}

We evaluate five models with clean features (Table~\ref{tab:models_overview}).

\begin{table}[t]
\centering
\caption{Model Overview}
\label{tab:models_overview}
\small
\begin{tabular}{lcccc}
\toprule
Model & ZS & Graph & Gene rep & Drug rep \\
\midrule
DistMult\textsuperscript{$\dagger$}  & No  & —      & Embed & Embed \\
TransE\textsuperscript{$\dagger$}    & No  & —      & Embed & Embed \\
R-GCN (bio)     & Yes & gene-mech & ESM-2 & Mol FP \\
Feature-MLP     & Yes & —         & ESM-2 & Mol FP \\
BioMolAMR       & Yes & gene-mech + drug-drug & ESM-2 & Mol FP \\
\bottomrule
\multicolumn{5}{l}{\small $\dagger$ Transductive: drug seen at train time as a lookup ID.}
\end{tabular}
\end{table}

\subsubsection{DistMult and TransE}
Standard transductive KG embedding baselines included for completeness.
At zero-shot evaluation time, these models score the held-out drug using its
trained embedding vector, which was optimised on seen classes. Their ZS scores
are therefore not true zero-shot and are marked accordingly.

\subsubsection{Feature-MLP}
A two-layer MLP with hidden dimension 256, GELU activation, and dropout 0.3.
Input is the concatenation of gene ESM-2 embedding and drug molecular
fingerprint; output is a scalar compatibility score. No graph structure is used.
Trained with margin ranking loss, 5 negatives per positive, full-batch
optimisation (batch size = 100,000; effectively one graph encoding per epoch).

\subsubsection{R-GCN (Bio Graph)}
A two-layer heterogeneous SAGE-based R-GCN that aggregates over gene-mechanism
and mechanism-gene edges. Drug embeddings are computed from molecular
fingerprints via a linear projection (no drug node neighbourhood aggregation).
Gene and drug embeddings are combined via bilinear scoring.

\subsubsection{BioMolAMR}
Our proposed model extends R-GCN with three additional components:
\begin{itemize}
  \item \textbf{Drug encoder}: graph SAGE over drug-drug chemical similarity
    edges (Tanimoto $> 0.20$), with molecular fingerprint initialisation.
  \item \textbf{Mechanism-bottleneck decoder}: a GATConv layer that routes
    gene-drug compatibility scores through mechanism representations, providing
    biological interpretability.
  \item \textbf{Multi-component loss}: ListNet ranking loss~\cite{cao2007listnet}
    + structural alignment loss (MSE to Tanimoto similarity matrix) +
    mechanism focal loss + prototype contrastive loss.
\end{itemize}

All inductive models use 5 seeds $\{42, 1, 2, 3, 4\}$, trained for up to
300 epochs with cosine annealing and early stopping (patience = 5 evaluations).

\section{Results}
\label{sec:results}

\subsection{Main Comparison}
\label{sec:main}

Table~\ref{tab:main} reports mean and standard deviation over 5 seeds.
Feature-MLP achieves the highest zero-shot performance on both metrics,
with ZS-all MRR = $0.069 \pm 0.006$---3.1$\times$ above the next-best inductive
model (R-GCN, $0.027$) and 3.1$\times$ above random ($0.022$).

\begin{table}[t]
\centering
\caption{Main Results (5 seeds). \textbf{Bold}: best inductive model.
$\dagger$: transductive (ZS not meaningful).}
\label{tab:main}
\begin{tabular}{lcccc}
\toprule
\textbf{Model} & \textbf{Test MRR} & \textbf{ZS-win}
  & \textbf{ZS-all}$^\ddagger$ & \textbf{Gen.\ Gap} \\
\midrule
\textbf{Feature-MLP}
  & \textbf{0.968}$\pm$0.002 & 0.137$\pm$0.005
  & \textbf{0.069}$\pm$0.006 & $-$0.899 \\
DistMult$^\dagger$
  & 0.767$\pm$0.007 & 0.141$\pm$0.023
  & 0.032$\pm$0.004 & $-$0.735 \\
R-GCN (bio)
  & 0.925$\pm$0.001 & 0.141$\pm$0.003
  & 0.027$\pm$0.001 & $-$0.898 \\
BioMolAMR
  & 0.912$\pm$0.011 & 0.125$\pm$0.005
  & 0.019$\pm$0.001 & $-$0.893 \\
TransE$^\dagger$
  & 0.887$\pm$0.008 & 0.120$\pm$0.038
  & 0.017$\pm$0.001 & $-$0.870 \\
\midrule
\textit{Random}
  & — & 0.083 & 0.022 & — \\
\midrule
\multicolumn{5}{l}{\textit{Prior work (leaky features, 5-way ZS-within)}} \\
Prior GNN~\cite{prior_work}
  & — & 0.775 & — & — \\
\bottomrule
\multicolumn{5}{l}{\small $\ddagger$ Primary metric: rank ZS drug among all 46.}
\end{tabular}
\end{table}

Several observations warrant attention:

\textbf{All inductive models show a large generalisation gap.} Test MRR
exceeds $0.90$ for graph models, but ZS-all MRR is near or below random for
BioMolAMR and TransE. The graph-based inductive models learn the training
graph structure well but fail to generalise. The Feature-MLP, despite weaker
training-set performance relative to graph models, generalises substantially
better.

\textbf{BioMolAMR underperforms random on ZS-all.} The proposed complex
architecture ($0.019$, below random $0.022$) loses to the simplest possible
inductive baseline. This is not a numerical accident: the standard deviation
across seeds is $\pm0.001$, confirming consistency of the failure.

\textbf{Transductive methods are not true zero-shot.} DistMult and TransE
assign non-zero scores to ZS drug classes because those classes have trained
lookup embeddings. Their ZS-all MRR ($0.032$ and $0.017$) reflects
extrapolation from the embedding space, not true inductive transfer.

\textbf{Prior work comparison.} Correcting for leakage, the best ZS-within MRR
drops from the prior-reported $0.775$ to $0.137$ (Feature-MLP), a reduction of
$5.7\times$. The corrected ZS-all MRR is $0.069$---3.1$\times$ above random,
indicating that genuine but modest zero-shot transfer is achievable with the
right feature representation.

\subsection{Per-Class Analysis}
\label{sec:perclass}

Fig.~\ref{fig:heatmap} shows per-class ZS-all MRR across all models.
Performance is highly heterogeneous: a small number of drug classes are well
predicted across all models, while the majority remain near random.

\textbf{Tractable classes.}
Lincosamide (Feature-MLP: $0.758$) and streptogramin A ($0.700$) are the
best-predicted classes. Both belong to the MLSB (macrolide-lincosamide-
streptogramin B) superfamily: their resistance genes overlap heavily with
trained macrolide resistance genes, and their molecular fingerprints have
moderate Tanimoto similarity to macrolide training drugs. Isoniazid-like
antibiotics ($0.370$) benefit from structural similarity to thioamide
antibiotics, a related training class.

\textbf{Hard classes.}
Fusidane ($\leq 0.12$ for all models), aminocoumarin ($\leq 0.15$), and
nitroimidazole ($\leq 0.10$) are near random for all methods. Fusidane has
only 11 resistance gene instances and targets elongation factor G, a mechanism
absent from most training drug classes. Aminocoumarin targets DNA gyrase B
(same target as fluoroquinolones) but via a structurally unrelated scaffold,
yielding very low Tanimoto similarity to any training drug.

\textbf{Chemical similarity as a predictor of transferability.}
Fig.~\ref{fig:tanimoto} shows the correlation between maximum Tanimoto
similarity of each ZS drug to the nearest training drug class and ZS-all MRR.
For Feature-MLP, Pearson $r = \texttt{[SEE FIGURE]}$, confirming that
chemical proximity is the primary driver of transferability.

\section{Ablation Studies}
\label{sec:ablations}

We evaluate three ablations of Feature-MLP (5 seeds each) to isolate the
source of zero-shot predictive power.

\begin{table}[t]
\centering
\caption{Feature Ablation Results (Feature-MLP variants, 5 seeds).
$\downarrow$ indicates drop from full model.}
\label{tab:ablations}
\begin{tabular}{lcc}
\toprule
\textbf{Variant} & \textbf{ZS-within MRR} & \textbf{ZS-all MRR} \\
\midrule
Feature-MLP (full)             & 0.137$\pm$0.004 & \textbf{0.069}$\pm$0.006 \\
FMLP: gene=0 (drug FP only)    & 0.140$\pm$0.003 & 0.020$\pm$0.002 $\downarrow$3.5$\times$ \\
FMLP: drug=0 (ESM-2 only)      & 0.010$\pm$0.000 & 0.010$\pm$0.000 $\downarrow$6.9$\times$ \\
FMLP: leaky gene features      & 0.159$\pm$0.030 & 0.112$\pm$0.000 $\uparrow$1.6$\times$ \\
\midrule
Random                         & 0.083           & 0.022           \\
\bottomrule
\end{tabular}
\end{table}

\subsection{Both Features Are Jointly Necessary}

\textbf{FMLP-gene-zero (drug fingerprints only).}
Zeroing ESM-2 gene embeddings causes ZS-all MRR to drop from $0.069$ to
$0.020 \approx $ random ($\downarrow$3.5$\times$), despite ZS-within MRR
remaining at $0.140$ (above the $0.083$ random baseline).
Drug fingerprints alone, without protein sequence context, are insufficient for
all-class zero-shot transfer. The drug fingerprint carries information about
\emph{which antibiotic class} is being queried, but without gene sequence
context the MLP cannot map it to the correct resistance genes.

\textbf{FMLP-drug-zero (ESM-2 only).}
Zeroing drug fingerprints causes ZS-all MRR to collapse to $0.010$ (below
random), with ZS-within MRR collapsing to $0.010$ as well (all seeds identical,
variance $= 0$). Without a drug-specific signal, the model assigns identical
scores to all genes regardless of the queried drug class, making ranking
effectively random with a fixed ordering bias.

The joint necessity of both features points to a \emph{pairwise interaction}
mechanism: the MLP learns gene-drug compatibility as a joint function of protein
sequence and chemical structure, neither of which suffices alone. This is
analogous to protein-ligand interaction prediction, where both binding site
structure and ligand chemistry must be represented.

\subsection{ESM-2 Gene Features Alone Insufficient (FMLP-drug-zero)}
As shown above, protein sequence information alone cannot generalise to novel
drug classes. ESM-2 embeddings encode evolutionary and functional protein
properties that are drug-agnostic by design---they are trained on sequences with
no knowledge of resistance phenotypes.

\subsection{Confirming Data Leakage (FMLP-leaky)}
Using the original 154-dimensional leaky gene features (containing the 46-bit
drug-class membership component), we train the Feature-MLP under identical
conditions. This directly quantifies the leakage effect: if ZS-all MRR is
substantially higher than $0.069$, the improvement is attributable entirely
to the leaked label information.

\subsection{Why Does Graph Structure Hurt?}
BioMolAMR and R-GCN underperform Feature-MLP despite richer architectures.
We hypothesise two mechanisms:

\textbf{Mechanism-channel smoothing.} Gene-mechanism-gene message passing
propagates information based on training-class resistance patterns. For seen
drug classes, this strengthens representations; for unseen classes, it blurs
drug-specific signal, averaging away the fingerprint information that Feature-MLP
exploits directly.

\textbf{Decoder overfit.} BioMolAMR's mechanism-bottleneck decoder learns
to route scores through 8 mechanism channels. During training, only 34 drug
classes are seen; the channel routing may overfit to their patterns and
produce miscalibrated scores for the 12 held-out classes.

These hypotheses are consistent with the observed behaviour: as model
complexity increases (Feature-MLP $\to$ R-GCN $\to$ BioMolAMR), ZS-all MRR
decreases monotonically despite training-set MRR increasing.

\section{Discussion}
\label{sec:discussion}

\subsection{Pairwise Interaction, Not Graph Topology}
Our ablations reveal that zero-shot AMR prediction requires \emph{both} gene
and drug feature modalities simultaneously: neither protein sequence nor
chemical structure alone is sufficient. This is consistent with the biological
reality of resistance: a gene's ability to resist a drug depends on the
specific biochemical interaction between the gene product (protein) and the
drug molecule. ESM-2 captures protein fold and functional context;
molecular fingerprints capture the chemical pharmacophore. Together, they
encode enough of this interaction to enable transfer to novel drug classes.

Our results contribute to a growing body of evidence that simple feature-based
models often match or outperform GNNs on biological link prediction tasks,
particularly in zero-shot settings~\cite{huang2020skipgnn,mao2023revisiting}.
The key insight for AMR is that the graph topology (which genes share
mechanisms, which drugs co-occur in resistance profiles) encodes training-class
structure that does not generalise to novel drug classes---and GNN message
passing propagates this non-transferable structure into node representations.

This suggests a principled design principle for zero-shot biological
prediction: \emph{invest in feature quality (pre-trained biological and chemical
encoders), not graph complexity}. For AMR specifically, replacing Morgan
fingerprints with learned molecular representations
(ChemBERTa~\cite{chithrananda2020chemberta}, GNN-based molecular encoders)
is a more promising direction than adding layers to the resistance graph GNN.

\subsection{Leakage as a Recurring Pitfall}
The leakage we identify is subtle but consequential: drug-class membership is
a biologically natural attribute of resistance genes, so including it in gene
features seems reasonable. However, in a zero-shot evaluation where drug
classes are the targets, any gene feature that encodes target drug-class
membership is a direct label leak. Analogous leakage can occur in:
\begin{itemize}
  \item Drug-protein interaction prediction when protein features encode known
    interaction partners,
  \item Disease gene prediction when gene features encode known disease
    associations, and
  \item Drug-side-effect prediction when drug features encode known
    side-effect profiles.
\end{itemize}
We recommend that any zero-shot biological prediction benchmark explicitly
audit whether features of the \emph{query entity} encode information about
the \emph{target entity}.

\subsection{Implications for AMR Surveillance}
From an applied perspective, our result is encouraging: screening novel
antibiotic candidates for potential resistance profiles requires only their
SMILES string (molecular structure) and a trained Feature-MLP. No gene
expression data, no genomic context, and no graph topology update is needed.
This is practically deployable at the early drug discovery stage.

\section{Conclusion}
\label{sec:conclusion}

We identified data leakage in prior zero-shot AMR prediction work,
introduced a rigorous clean benchmark, and showed that a simple
Feature-MLP achieves state-of-the-art zero-shot performance
($\text{ZS-all MRR} = 0.069$), outperforming all tested GNNs by $3.6\times$.
Ablations establish that \emph{both} ESM-2 gene embeddings and molecular drug
fingerprints are jointly necessary; neither alone achieves above-random ZS-all
MRR. Graph topology provides no additional zero-shot benefit and can harm
generalisation. Feature leakage in prior work inflated ZS MRR by $5.7\times$
(from $0.137$ to the reported $0.775$). The ``right tool'' for zero-shot AMR
is pairwise interaction modelling between protein sequence and chemical
structure---a direction better served by improved biological and molecular
encoders than by more complex graph architectures.

\section*{Acknowledgments}
CARD data provided by the McMaster Resistance Research Group under
Creative Commons Attribution 4.0 Licence.

\bibliographystyle{IEEEtran}
\bibliography{references}

\end{document}
